\section{Related Work}\label{sec:related_work}

Financial time-series generators are commonly evaluated by examining whether their outputs reproduce established empirical properties of asset returns.
\citet{takahashi2019modeling} assess generated series against stylized facts including heavy tails, volatility clustering, leverage effects, and gain/loss asymmetry.
Similarly, \citet{wiese2020quant} evaluate both marginal and temporal properties and show that their Quant GAN can reproduce volatility clustering, leverage effects, and serial dependence.
More recently, \citet{kwon2024stylized} systematically compare GAN architectures and find that their ability to recover univariate and multivariate stylized facts depends strongly on the generator architecture.
Diffusion-based approaches have also been evaluated in this manner: \citet{tanaka2025cofindiff} examine whether generated series preserve financial properties under specified trend and volatility conditions, while \citet{zhang2026sfag} incorporate selected stylized facts directly into the training objective.
These studies establish the importance of financially grounded evaluation, but the resulting diagnostics are generally presented as plots, summary statistics, tests of presence, or uncalibrated discrepancies.
They therefore provide limited guidance on whether an observed synthetic--real difference exceeds the variation naturally arising between finite samples of real data.

Several studies instead combine multiple evaluation criteria.
\citet{ericson2024deep} introduce financial modelling indicators covering distributional distance, autocorrelation structure, and backtesting, and compare several conditional generators on historical yield-curve data and simulated processes.
\citet{hounwanou2025evaluating} similarly combine maximum mean discrepancy, autocorrelation and volatility-clustering diagnostics, and downstream evaluations based on portfolio optimisation and volatility forecasting.
These approaches broaden evaluation beyond any single measure of realism, but their components retain different numerical scales and interpretations.
Consequently, they produce a collection of complementary indicators rather than a common reference against which the severity of each discrepancy can be assessed.

Our framework builds on this multi-criteria perspective but differs in both its diagnostic focus and its reference for interpretation.
Rather than reporting heterogeneous financial diagnostics separately or
collapsing them into a single global measure, it evaluates distinct financial
properties individually and benchmarks each synthetic--real discrepancy
against the variation observed between matched resamples of the real ticker
corpus.
This common empirical reference gives the resulting metrics a consistent
interpretation and indicates which targeted statistical signatures exhibit
larger or smaller discrepancies.
The framework therefore provides a calibrated diagnostic assessment of generator fidelity, rather than an unreferenced collection of indicators or a single aggregate verdict.
